\section{Introducción}

Los puntos cuánticos de carbono (CQDs, por sus siglas en inglés) son nanoestructuras de carbono cerodimensionales con tamaños inferiores a $10$~nm que han despertado un enorme interés en la física de materiales y la nanotecnología. A diferencia de los puntos cuánticos semiconductores tradicionales basados en metales pesados tóxicos (como CdSe o PbS), los CQDs ofrecen alta biocompatibilidad, baja toxicidad, estabilidad química y una fotoluminiscencia sintonizable. Estas propiedades los convierten en candidatos prometedores para aplicaciones en optoelectrónica, bioimagen, senosado químico y catálisis.
La incorporación de heteroátomos como nitrógeno ($N$) y azufre ($S$) en la matriz carbonosa (dopaje y co-dopado) permite modificar la densidad de carga local, alterar la brecha de banda óptica y generar estados de trampa superficiales que mejoran la eficiencia cuántica de emisión. Entre los métodos de síntesis bottom-up, la pirólisis termoquímica asistida por microondas destaca por su rapidez, bajo consumo energético y capacidad para promover reacciones de condensación y carbonización en cuestión de minutos.
Frente al calentamiento hidrotermal convencional, la pirólisis asistida por microondas proporciona perfiles de temperatura homogéneos y ultra-rápidos que aceleran la deshidratación y la carbonización del núcleo $sp^2$ en escala de minutos, permitiendo además el aprovechamiento de biomasa y bioprecursores sostenibles \cite{ren2024carbon}.


En este trabajo se presenta un estudio experimental comparativo de la síntesis y caracterización de puntos cuánticos de carbono utilizando tres rutas termoquímicas distintas:
\begin{enumerate}
    \item \textbf{Ruta Base (Ácido Cítrico y Urea):} Evaluación de la relación molar (1:6 y 1:10) para analizar la transición entre matrices masivas monolíticas afectadas por auto-extinción por agregación (ACQ) y soluciones coloidales fluorescentes por confinamiento estérico.
    \item \textbf{Ruta Verde (Ácido Cítrico y Clara de Huevo):} Empleo de la ovoalbúmina como bioprecursor sostenible de nitrógeno y azufre (N,S-CQDs), determinando la ventana cinética óptima (\textit{sweet spot}) a 1:15~minutos de pirólisis y su \textit{bandgap} óptico ($E_g = 3{,}334$~eV).
    \item \textbf{Ruta Ácida (Ácido Cítrico, Urea y Vinagre):} Análisis del efecto del medio ácido en la modificación de los estados de superficie y el corrimiento batocrómico resultante ($+57$~nm en UV-Vis).
\end{enumerate}
Asimismo, ante la ausencia de un espectrofluorómetro comercial, se implementó y validó una plataforma metrológica portátil de bajo costo basada en \textbf{Fotometría Digital RGB en Cámara Oscura}. Mediante la sustracción de fondo ($Y_{\text{neta}}$) y el análisis del ratio ratiométrico verde-azul ($G/B$), se cuantificó tanto la cinética de reacción como el desplazamiento espectral entre las diferentes rutas sintéticas, ofreciendo una metodología reproducible respaldada por la norma internacional ITU-R BT.601 \cite{poynton2012digital} y la literatura analítica actual.